% -------------------------------------------------------------- %
%      csd5748 - HY-225 - Spring 2025 - CSD-UOC - Ex. 08-09      %
% -------------------------------------------------------------- %

\documentclass{article}

\usepackage[utf8]{inputenc} 
\usepackage[english, greek]{babel}
\usepackage{subcaption}
\usepackage{hyperref}
\usepackage{xcolor}
\hypersetup{ hidelinks, colorlinks=true, linkcolor=blue, urlcolor=blue }
\captionsetup[subfigure]{labelformat=empty}
\usepackage{amsmath}
\usepackage{amsfonts}
\usepackage{amssymb}
\usepackage{array}
\usepackage{graphicx}
\usepackage[a4paper, left=1.1cm, right=1.1cm, top=2cm, bottom=2cm]{geometry}
\usepackage[absolute, overlay]{textpos}
\usepackage{placeins}
\usepackage{listings}
\usepackage{xcolor}
\setlength{\jot}{0pt}
\setlength{\mathsurround}{0.1cm}
\setlength{\jot}{0.1cm}
\newcommand{\skato}{\hspace{0.1cm}}
\newcommand{\skata}{\vspace{1cm}}

\begin{document}

% -------------------------------------------------------------- %

\begin{center}
\setlength{\fboxrule}{2pt}
\setlength{\fboxsep}{15pt} 
\fbox
{
      {\noindent {\Huge \textbf{\textlatin{HY225 -} Ασκήσεις 8 και 9 \textlatin{- csd5748}}}}
}
\end{center}
\vspace{1.100cm}

% -------------------------------------------------------------- %

\begin{center}
\setlength{\fboxrule}{1.1pt}
\setlength{\fboxsep}{7pt}
    \fbox{{\textbf{Άσκηση 8.}}}
    \vspace{0.75cm}
\end{center}

{\underline{\textbf{Άσκηση 8.5α:}}}
\vspace{1.25cm}

{
\Huge
\begin{table}[h!]
\centering
\begin{tabular}{|c|c|c|c|c|c|c|c|}
\hline
\textlatin{Opcode} & \textlatin{aluSrc} & \textlatin{rfSrc} & \textlatin{mRd} & \textlatin{mWr} & \textlatin{rfWr} & \textlatin{brOp} & \textlatin{aluOp} \\
\hline
0000011 & 1 & 1 & 1 & 0 & 1 & 0 & 000 \\
0100011 & 1 & \textlatin{x} & 0 & 1 & 0 & 0 & 000 \\
1100011 & 0 & \textlatin{x} & 0 & 0 & 0 & 1 & 001 \\
0010011 & 1 & 0 & 0 & 0 & 1 & 0 & 010 \\
0110011 & 0 & 0 & 0 & 0 & 1 & 0 & 110 \\
\textlatin{All others.} & \textlatin{x} & \textlatin{x} & 0 & 0 & 0 & 0 & 000 \\
\hline
\end{tabular}
\caption*{Πίνακας Αληθείας: \textlatin{Main Ctrl}.}
\end{table}
}

% -------------------------------------------------------------- %

\vspace{0.321cm}

\setlength{\fboxrule}{1pt}
\begin{figure}[h]
  \centering
  \fbox{\includegraphics[width=0.75\textwidth]{8.5a.jpg}}
  \caption*{Κύκλωμα: \textlatin{Main Ctrl} - σχεδιασμένο στο \href{https://circuitverse.org/simulator}{\textlatin{CircuitVerse}}.}
\end{figure}

% -------------------------------------------------------------- %

\newpage
\textbf{   }
\vspace{1cm}
\underline{\textbf{Άσκηση 8.5b:}}
\vspace{1.25cm}


\begin{table}[h!]
\centering
\begin{tabular}{|c|c|c|c|c|}
\hline
\textlatin{brOp} & \textlatin{zero} & \textlatin{sign} & \textlatin{funct3} & \textlatin{pcSrc} \\
\hline
0 & \textlatin{x} & \textlatin{x} & \textlatin{xxx} & 0 \\
1 & 1 & \textlatin{x} & 000 & 1 \\
1 & 0 & \textlatin{x} & 000 & 0 \\
1 & 0 & \textlatin{x} & 001 & 1 \\
1 & 1 & \textlatin{x} & 001 & 0 \\
1 & \textlatin{x} & 1 & 10\textlatin{x} & 1 \\
1 & \textlatin{x} & 0 & 10\textlatin{x} & 0 \\
1 & \textlatin{x} & 0 & 101 & 1 \\
1 & \textlatin{x} & 1 & 101 & 0 \\
\hline
\end{tabular}
\caption*{Πίνακας Αληθείας Διακλαδώσεων.}
\end{table}

% -------------------------------------------------------------- %

\vspace{0.321cm}

\setlength{\fboxrule}{1pt}
\begin{figure}[h]
  \centering
  \fbox{\includegraphics[width=0.55\textwidth]{8.5b.jpg}}
  \caption*{Κύκλωμα Σήματος Διακλαδώσεων - σχεδιασμένο στο \href{https://circuitverse.org/simulator}{\textlatin{CircuitVerse}}.}
\end{figure}

% -------------------------------------------------------------- %

\newpage
\underline{\textbf{Άσκηση 8.5g:}}
\vspace{1.25cm}


\begin{table}[h!]
\centering
\begin{tabular}{|c|c|c|c|}
\hline
\textlatin{aluOp} & \textlatin{funct3} & \textlatin{funct7} & \textlatin{aluMd} \\
\hline
000 & \textlatin{xxx} & \textlatin{xxxxxxxx} & 00\textlatin{x}0 \\
001 & \textlatin{xxx} & \textlatin{xxxxxxxx} & 00\textlatin{x}1 \\
010 & 000 & \textlatin{xxxxxxxx} & 00\textlatin{x}0 \\
010 & 110 & \textlatin{xxxxxxxx} & 010\textlatin{x} \\
010 & 111 & \textlatin{xxxxxxxx} & 011\textlatin{x} \\
010 & \textlatin{other} & \textlatin{xxxxxxxx} & 00\textlatin{x}0 \\
110 & 000 & 0000000  & 00\textlatin{x}0 \\
110 & 000 & 0100000  & 00\textlatin{x}1 \\
110 & 000 & 0000001  & 1\textlatin{xxx} \\
110 & 000 & \textlatin{other} & 00\textlatin{x}0 \\
110 & 110 & \textlatin{xxxxxxxx} & 010\textlatin{x} \\
110 & 111 & \textlatin{xxxxxxxx} & 011\textlatin{x} \\
110 & \textlatin{other} & \textlatin{xxxxxxxx} & 00\textlatin{x}0 \\
\hline
\end{tabular}
\caption*{Πίνακας Αληθείας \textlatin{ALU Ctrl}.}
\end{table}

% -------------------------------------------------------------- %

\vspace{0.5cm}

\setlength{\fboxrule}{1pt}
\begin{figure}[h]
  \centering
  \fbox{\includegraphics[width=0.55\textwidth]{8.5g.jpg}}
  \caption*{Κύκλωμα: \textlatin{ALU Ctrl} - σχεδιασμένο στο \href{https://circuitverse.org/simulator}{\textlatin{CircuitVerse}}.}
\end{figure}

% -------------------------------------------------------------- %

\newpage
\underline{\textbf{Άσκηση 8.6:}}
\vspace{1.350cm}


\setlength{\fboxrule}{1pt}
\begin{figure}[h]
  \centering
  \fbox{\includegraphics[width=0.75\textwidth]{8.6.jpg}}
  \captionsetup{justification=centering}
  \caption*{[ Εικόνα από \href{https://www.csd.uoc.gr/~hy225/23a/08a_scpi.pdf\#page=2}{\underline{εδώ}}. ]}
\end{figure}

\vspace{1.0cm}
\textbf{Λοιπόν, τώρα που προσθέθηκαν δύο νέοι πολυπλέκτες στο κύκλωμα (πάνω-αριστερά και κάτω-κέντρο), θα πρέπει να προσθέσουμε και δύο νέα σήματα για τους ελέγχους των πολυπλεκτών (\textlatin{jal} και \textlatin{jalr}), καθώς και να τροποποιηθούν κάποια άλλα (\textlatin{rfWr} και \textlatin{aluSrc}), ώστε να λειτουργούν όλα όπως πρέπει. Έτσι, το νέο σήμα \textlatin{jal} επιλέγει πια τιμη θα αποθηκευτεί στον καταχωρητή \textlatin{rd}, η έξοδος της \textlatin{ALU} ή η τιμή \textlatin{PC+4}. Από την άλλη, το νέο σήμα ελέγχου \textlatin{jalr} επιλέγει την επόμενη διεύθυνση \textlatin{PC}, μεταξύ της εξόδου της \textlatin{ALU} και του \textlatin{PC+4}. Τέλος, το υπάρχον σήμα ελέγχου \textlatin{rfWr} πρέπει να ενεργοποιείται και για τις δύο εντολές \textlatin{jal} και \textlatin{jalr} κατάλληλα.}


% -------------------------------------------------------------- %

\newpage
\underline{\textbf{Άσκηση 8.7:}}
\vspace{1.35cm}


\textbf{Έχουμε ως Δεδομένα Καθυστερήσεων:}
\vspace{0.1cm}

\textlatin{- Instruction Memory: 500 ps.}
\vspace{0.1cm}

\textlatin{- Data Memory: 500 ps.}
\vspace{0.1cm}

\textlatin{- Register File (read): 300 ps.}
\vspace{0.1cm}

\textlatin{- Register File (write): 200 ps.}
\vspace{0.1cm}

\textlatin{- ALU: 400 ps.}
\vspace{0.1cm}

\textlatin{- Adder (PC): 400 ps.}
\vspace{0.1cm}

\textlatin{- Control Logic: 200 ps.}
\vspace{0.1cm}

\textlatin{- Multiplexers - PC: 0 ps} \textbf{ (αγνοούνται).}
\vspace{0.75cm}

\textbf{\textlatin{a}) Η εντολή: \textlatin{500 ps.}}

\vspace{0.1cm}
\textbf{\textlatin{b}) Τα σήματα ελέγχου: \textlatin{500 + 200 = 700 ps.}}

\vspace{0.1cm}
\textbf{\textlatin{c}) Οι τιμές των καταχωρητών πηγής: \textlatin{500 + 300 = 800 ps.}}

\vspace{0.1cm}
\textbf{\textlatin{d}) Η έξοδος της \textlatin{ALU}: \textlatin{500 + 300 + 400 = 1200 ps.}}

\vspace{0.1cm}
\textbf{\textlatin{e}) Οι τιμές \textlatin{PC+4} και \textlatin{PC+2×Imm12}: \textlatin{400 ps} και \textlatin{500 + 400 = 900 ps} αντίστοιχα.}

\vspace{0.1cm}
\textbf{\textlatin{f}) Η τιμή γιά τον επόμενο \textlatin{PC}: \textlatin{500 + 300 + 400 = 1200 ps.}}

\vspace{0.1cm}
\textbf{\textlatin{g}) Η έξοδος της \textlatin{Data Memory}: \textlatin{500 + 300 + 400 + 500 = 1700 ps.}}

\vspace{0.1cm}
\textbf{\textlatin{h}) Η είσοδος \textlatin{Write Data} του Αρχείου Καταχωρητών: \textlatin{500 + 300 + 400 + 500 = 1700 ps.}}

\vspace{0.1cm}
\textbf{\textlatin{i}) Ολοκλήρωση τυχόν εγγραφής στο Αρχείο Κατ.: \textlatin{500 + 300 + 400 + 500 + 200 = 1900 ps.}}

\vspace{0.75cm}
\textbf{- Χειρότερη Περίπτωση: \textlatin{1900 ps.}}

\vspace{0.1cm}
\textbf{- Ελάχιστη περίοδος ρολογιού: \textlatin{1.9 ns.}}

\vspace{0.1cm}
\textbf{- Μέγιστη Συχνότητα Ρολογιού: } \textlatin{f = 1 / 1.9 = 0.526 = 526 MHz}

\vspace{0.1cm}
\textbf{- Με Εντολή Διαίρεσης: \textlatin{500 + 300 + 4400 + 500 + 200 = 5900 ps = 5.9 ns.}}

\vspace{0.1cm}
\textbf{- Νέα συχνότητα: } \textlatin{f = 1 / 5.9 = 0.169 = 169 MHz}

\vspace{0.1cm}
\textbf{- Τελικά, πιο αργό κατά: 526 / 169 = 3.112 φορές.}

\vspace{10.1cm}
\hline

% -------------------------------------------------------------- %

\newpage

\begin{center}
\setlength{\fboxrule}{1.1pt}
\setlength{\fboxsep}{7pt}
    \fbox{{\textbf{Άσκηση 9.}}}
    \vspace{0.75cm}
\end{center}

\underline{\textbf{Άσκηση 9.9:}}
\vspace{1.75cm}


\textbf{(1)\hspace{0.1cm} Οι βαθμίδες (στάδια) της \textlatin{pipeline} του επεξεργαστή (από τον \href{https://ripes.me/}{\textlatin{Ripes}}) είναι: }

\vspace{0.5cm}

\textbf{α. \textlatin{Instruction Fetch (IF)}. Λειτουργία: Ανακτά την εντολή από την \textlatin{Instruction Memory (IM)} χρησιμοποιώντας τη διεύθυνση του \textlatin{Program Counter (PC)} και ενημερώνει τον \textlatin{PC} κατά \textlatin{(PC+4)} για την επόμενη εντολή και, τέλος, περνά την εντολή στο επόμενο στάδιο.}

\vspace{0.25cm}

\textbf{β. \textlatin{Instruction Decode (ID)}. Λειτουργία: Αποκωδικοποιεί την εντολή \textlatin{(opcode, registers, immediates)} μέσω του \textlatin{Decoder} καθώς και διαβάζει τους καταχωρητές πηγής \textlatin{(reg1, reg2)} από το \textlatin{Register File (RF)} και περνά τα δεδομένα στο επόμενο στάδιο.}

\vspace{0.25cm}

\textbf{γ. \textlatin{Execute (EX)}. Λειτουργία: Εκτελεί αριθμητικές ή λογικές πράξεις στην \textlatin{ALU} με βάση τους \textlatin{(Op1, Op2)} αλλά και υπολογίζει διευθύνσεις για \textlatin{load/store} ή ελέγχει διακλαδώσεις \textlatin{(Branch)}, περνώντας το αποτέλεσμα στο επόμενο στάδιο.}

\vspace{0.25cm}

\textbf{δ. \textlatin{Memory Access (MEM)}. Λειτουργία: Πραγματοποιεί πρόσβαση στη \textlatin{Data Memory (DM)} για εντολές \textlatin{load} (ανάγνωση) ή \textlatin{store} (εγγραφή), εαν χρειάζεται, και περνά τα δεδομένα ή τα αποτελέσματα στο επόμενο στάδιο.}

\vspace{0.25cm}

\textbf{ε. \textlatin{Write Back (WB)}. Λειτουργία: Εγγράφει το αποτέλεσμα από την \textlatin{ALU} ή την \textlatin{DM} πίσω στο \textlatin{Register File}. Αυτό είναι και το τελευταίο στάδιο (βαθμίδα) της \textlatin{pipeline} του συγκεκριμένου επεξεργαστή.}


\vspace{2cm}


\textbf{(2)\hspace{0.1cm} Οι βαθμίδες (στάδια) της \textlatin{pipeline} που είναι απαραίτητες για ολες τις εντολές είναι:}

\vspace{0.5cm}

\textbf{α. \textlatin{Instruction Fetch (IF)}, αφού ολες οι εντολές πρέπει να ανακτηθούν από τη μνήμη \textlatin{(Instruction Memory)} για να ξεκινήσει η εκτέλεσή τους.}

\vspace{0.25cm}

\textbf{β. \textlatin{Instruction Decode (ID)}, αφού η αποκωδικοποίηση είναι απαραίτητη για να κατανοηθεί τι κάνει η εντολή και να διαβαστούν οι καταχωρητές, για να προετοιμαστούν για εκτέλεση.}

\vspace{0.25cm}

\textbf{γ. \textlatin{Execute (EX)}, αφού η \textlatin{ALU} χρησιμοποιείται για όλες τις εντολές και είναι έτσι πάντα ενεργή.}

 
\vspace{2cm}


\textbf{(3\textlatin{i}) Οι βαθμίδες (στάδια) της \textlatin{pipeline} που δεν χρειάζονται για εντολές αριθμητικών πράξεων:}

\vspace{0.5cm}

\textbf{α. \textlatin{Memory Access (MEM)}, αφού οι αριθμητικές πράξεις δεν απαιτούν πρόσβαση στη μνήμη \textlatin{(Data Memory)}. Το αποτέλεσμα της \textlatin{ALU} απλώς περνά στο επόμενο στάδιο.}


\vspace{1cm}


\textbf{(3\textlatin{ii})\hspace{0.1cm} Οι βαθμίδες (στάδια) της \textlatin{pipeline} που δεν χρειάζονται για εντολές \textlatin{store} είναι:}

\vspace{0.5cm}

\textbf{α. \textlatin{Write Back (WB)}, αφού οι εντολές αυτές δεν δημιουργούν αποτέλεσμα που να πρέπει να γραφτεί σε κάποιο καταχωρητή.}


% -------------------------------------------------------------- %


\newpage
\textbf{   }
\vspace{1cm}
\textbf{(4)\hspace{0.1cm} Οι αλληλουχίες εντολών χωρίς αλληλεξαρτήσεις και οι εκτελέσεις τους στον \href{https://ripes.me/}{\textlatin{Ripes}:}}
\vspace{1cm}
\begin{figure}[h]
  \centering
  \setlength{\fboxrule}{1pt}
  \fbox{\includegraphics[width=0.63\textwidth]{9.9da.jpg}} \\
  \vspace{0.1cm}
  \fbox{\includegraphics[width=0.63\textwidth]{9.9db.jpg}} \\
  \vspace{0.1cm}
  \fbox{\includegraphics[width=0.63\textwidth]{9.9dc.jpg}}
\end{figure}


% -------------------------------------------------------------- %


\newpage
\textbf{   }
\vspace{1cm}
\textbf{(5)\hspace{0.1cm} Οι νέες εξαρτημένες εντολές του ερ. (4α) και οι εκτελέσεις τους στον \href{https://ripes.me/}{\textlatin{Ripes}} (μαζί):}

\setlength{\fboxrule}{1pt}
\begin{figure}[h]
  \centering
  \fbox{\includegraphics[width=0.75\textwidth]{9.9e.jpg}}
\end{figure}


% -------------------------------------------------------------- %


\newpage
\textbf{(6)\hspace{0.1cm} Οι εντολές πριν και μετά την αναδιάταξη μαζί με τις εκτελέσεις τους στον \href{https://ripes.me/}{\textlatin{Ripes}} είναι:}
\vspace{1cm}

\begin{figure}[h]
\hspace{-0.5cm}
  \centering
  \begin{minipage}[b]{0.387\textwidth}
    \centering
    \fbox{\includegraphics[width=\linewidth]{9.9fa0.jpg}}
  \end{minipage}
  \hspace{0.3cm}
  \begin{minipage}[b]{0.58\textwidth}
    \centering
    \fbox{\includegraphics[width=\linewidth]{9.9fa1.jpg}}
  \end{minipage}
  
  \vspace{0.3cm}
  
\hspace{-0.5cm}
  \centering
  \begin{minipage}[b]{0.39\textwidth}
    \centering
    \fbox{\includegraphics[width=\linewidth]{9.9fb0.jpg}}
  \end{minipage}
  \hspace{0.3cm}
  \begin{minipage}[b]{0.58\textwidth}
    \centering
    \fbox{\includegraphics[width=\linewidth]{9.9fb1.jpg}}
  \end{minipage}
\end{figure}


% -------------------------------------------------------------- %


\vspace{1.750cm}

\textbf{(7) \hspace{0.1cm} Οι περισσότεροι από τους 189 - [5 + (125 - 1)] - 6 = 54 χαμένους κύκλους οφείλονται σε προβλήματα των 16 διακλαδώσεων στο βρόχο αλλά και στις κλήσεις συναρτήσεων που χρησιμοποιούνται στο \textlatin{factorial.s}. Μία, λοιπόν, τεχνική για να βελτιωθεί αυτό είναι η επόμενη εντολή μετά τη διακλάδωση να εκτελείται πάντα. Με αυτόν τον τρόπο, θα γλιτώναμε περίπου 16 κύκλους.}


\vspace{7.1cm}
\hline


% -------------------------------------------------------------- %


\newpage

\begin{center}
    \vspace*{\fill}
    \setlength{\fboxrule}{1.1pt}
    \setlength{\fboxsep}{7pt}
    \fbox{\textbf{\LARGE{ΤΕΛΟΣ}}}
    \vspace*{\fill}
\end{center}


% -------------------------------------------------------------- %

\end{document}

% -------------------------------------------------------------- %
%                           The End :)                           %
% -------------------------------------------------------------- %